% !TeX root = ../../main.tex

Die einzelnen Experimente verfolgen unterschiedliche Schwerpunkte.
Wie in Abschnitt~\ref{subsec:experimentFunctionsImpl} beschrieben, untersucht das erste Experiment die Fehlerklassen,
ohne den Einsatz eines Retry-Mechanismus.
Das dient als Baseline, mit der die Ergebnisse verglichen werden.
Das zweite Experiment(\texttt{experiment\_1}) untersucht den grundlegenden Einfluss von Retries auf die Transaktionen.
Dadurch kann bestimmt werden, wie sich das Verhalten eines Retries ohne Verzögerungsmechaniken auf das System auswirkt.
Im dritten Experiment wird der Einfluss von Delay ermittelt, um festzustellen, wie unterschiedliche Wartezeiten sich
auf Systemverhalten auswirken.
Das letzte Experiment vergleicht die einzelnen Retry-Strategien miteinander.
Dadurch kann in der anschließenden Analyse untersucht werden, ob sich bestimmte Strategien unter den gewählten
Bedingungen als vorteilhafter oder nachteiliger erweisen.

Für die Durchführung der Experimente stehen die in Codeausschnitt~\ref{lst:params} dargestellten Parameter zur
Verfügung.
Aus deren möglichen Werten werden die einzelnen Konfigurationen für die Experimente abgeleitet.
Für das Baseline-Experiment werden ausschließlich die verschiedenen Concurrency-Werte betrachtet.
Bei fünf Concurrency-Werten und fünf Wiederholungen ergeben sich somit insgesamt 25 Konfigurationen.
Im Experiment zur Untersuchung des Retry-Einflusses wird zusätzlich die maximale Anzahl der möglichen Retries variiert.
Dabei werden die Werte 1, 3 und 5 für jede der fünf Concurrency-Stufen untersucht.
Der Wert 0 wird nicht erneut betrachtet, da dieser bereits durch das Baseline-Experiment abgedeckt wird und somit
keine zusätzlichen Erkenntnisse liefert.
Daraus ergeben sich $(5 \cdot 3 \cdot 5 = 75)$ Konfigurationen.
Für die Untersuchung des Delays werden die vier von null verschiedenen Werte des initialen Delays für jede
Concurrency-Stufe getestet.
Auch hier wird jede Konfiguration fünfmal wiederholt, wodurch sich insgesamt $(5 \cdot 4 \cdot 5 = 100)$
Konfigurationen ergeben.
Im letzten Experiment werden erneut alle fünf Concurrency-Werte betrachtet.
Zusätzlich werden zwei Retry-Strategien sowie die drei von null verschiedenen Retry-Counts untersucht.
Daraus ergeben sich $(5 \cdot 2 \cdot 3 \cdot 5 = 150)$ Konfigurationen.
Aus den 350 Konfigurationen und den jeweils 500 ausgeführten Transaktionen entstehen insgesamt 175.000
Transaktionsausführungen, deren Ergebnisse in den Logdateien erfasst und anschließend ausgewertet werden.

Für die Auswertung werden die in Abschnitt~\ref{sec:purpose} beschriebenen Kenngrößen verwendet.
Dazu gehören die Erfolgsrate, die Retry-Verteilung, die durchschnittliche Ausführungszeit, die Quartile der
Ausführungszeit, der durch Retries entstehende Overhead sowie der Durchsatz pro Experiment.
Die durchschnittliche Ausführungszeit, die Quartile und der Overhead werden in Millisekunden angegeben.
Die Erfolgsrate wird als Prozentwert dargestellt, während die Retry-Verteilung die Anzahl der Transaktionen je
durchgeführter Retries angibt.
Der Durchsatz wird als Anzahl erfolgreicher Transaktionen pro Sekunde und Experiment angegeben.
